% GENERATED FILE. Edit canonical lessons, then run npm run generate:books.
% canonical-source-hash: fnv1a64:17ea827e0bfa8cae
% canonical-lessons: JA-C317-sansei-suru, JA-C317-hantai-suru, JA-C317-giron, JA-C317-teian-suru, JA-C317-ketsuron

\chapter{Agreeing and Disagreeing}
\label{ch:ja-agreeing-and-disagreeing}
\hlchaptermodality{317}

\begin{chapteropening}
\textbf{By the end of this chapter:} \emph{I can say whether I agree with an idea, and why.}

Complete the last lesson of chapter 317: I can say whether I agree with an idea, and why.
\end{chapteropening}

\section[\texorpdfstring{\ja{さんせいする} (sansei suru)}{さんせいする (sansei suru)}]{\ja{さんせいする} (sansei suru) — to agree; to be in favour}

\label{lesson:JA-C317-sansei-suru}



\begin{quote}
Before the new one: say the Japanese for serious, then the Japanese for selfish.
\end{quote}

\subsection*{You'll want to know: \ja{さんせいする}}
\textbf{\ja{さんせいする}} — \emph{sansei suru} — \textquotedblleft{}to agree; to be in favour\textquotedblright{}.

\begin{quote}
\textbf{\ja{あにの} \ja{いけんは} \ja{いいと} \ja{おもったので、さんせいしました。}} — \emph{ani no iken wa ii to omotta node, sansei shimashita} — \textquotedblleft{}I thought my brother's idea was good, so I agreed with it.\textquotedblright{}
\end{quote}

Polite \textbf{\ja{さんせいします}}, \emph{sanseishimasu}; a \ja{する} verb. Past (polite): \textbf{\ja{ははは} \ja{さんせいしました。}}, \emph{haha wa sansei shimashita}.

\begin{cousinweb}
\textbf{\ja{さんせいする}} (\emph{sansei suru}) is built from a Sino-Japanese noun and the native verb \emph{suru}, \textquotedblleft{}to do\textquotedblright{}. The noun, \emph{sansei}, is two Chinese characters read with their on'yomi, readings borrowed from Chinese: \emph{san}, \textquotedblleft{}approve\textquotedblright{} (Middle Chinese \emph{tsanH}) and \emph{sei}, \textquotedblleft{}become\textquotedblright{} (\emph{dzyeng}) (KANJIDIC2; JMdict; Middle Chinese in Baxter's transcription).
\end{cousinweb}

\begin{grammarlens}[title={What you've built}]
One more word for giving reasons and opinions.
\end{grammarlens}

\subsection*{Your turn}
\begin{itemize}
\raggedright
  \item \emph{Say it:} \emph{sansei suru}
  \item \emph{Say it:} \emph{sansei suru}, once more
  \item \emph{Say it:} \emph{ani no iken wa ii to omotta node, sansei shimashita} — I thought my brother's idea was good, so I agreed with it.
  \item \emph{Recall:} say the Japanese for serious, then the Japanese for selfish, then say \emph{sansei suru} again
\end{itemize}

\begin{tcolorbox}[breakable,colback=teal!4,colframe=teal!35!black,title={Before you move on}]
What does \ja{さんせいする} mean? (\textquotedblleft{}To agree; to be in favour\textquotedblright{}.) Say it once more.
\end{tcolorbox}
\section[\texorpdfstring{\ja{はんたいする} (hantai suru)}{はんたいする (hantai suru)}]{\ja{はんたいする} (hantai suru) — to be against; to oppose}

\label{lesson:JA-C317-hantai-suru}



\begin{quote}
Before the new one: say the Japanese for selfish, then the Japanese for to agree.
\end{quote}

\subsection*{You'll want to know: \ja{はんたいする}}
\textbf{\ja{はんたいする}} — \emph{hantai suru} — \textquotedblleft{}to be against; to oppose\textquotedblright{}.

\begin{quote}
\textbf{\ja{あぶないから、ははは} \ja{はんたいしました。}} — \emph{abunai kara, haha wa hantai shimashita} — \textquotedblleft{}My mother was against it because it is dangerous.\textquotedblright{}
\end{quote}

Polite \textbf{\ja{はんたいします}}, \emph{hantaishimasu}; a \ja{する} verb. Past (polite): \textbf{\ja{あには} \ja{はんたいしました。}}, \emph{ani wa hantai shimashita}.

\begin{cousinweb}
\textbf{\ja{はんたいする}} (\emph{hantai suru}) is built from a Sino-Japanese noun and the native verb \emph{suru}, \textquotedblleft{}to do\textquotedblright{}. The noun, \emph{hantai}, is two Chinese characters read with their on'yomi, readings borrowed from Chinese: \emph{han}, \textquotedblleft{}anti-\textquotedblright{} (Middle Chinese \emph{pjonX}) and \emph{tai}, \textquotedblleft{}opposite\textquotedblright{} (\emph{twojH}) (KANJIDIC2; JMdict; Middle Chinese in Baxter's transcription).
\end{cousinweb}

\begin{grammarlens}[title={What you've built}]
One more word for giving reasons and opinions.
\end{grammarlens}

\subsection*{Your turn}
\begin{itemize}
\raggedright
  \item \emph{Say it:} \emph{hantai suru}
  \item \emph{Say it:} \emph{hantai suru}, once more
  \item \emph{Say it:} \emph{abunai kara, haha wa hantai shimashita} — My mother was against it because it is dangerous.
  \item \emph{Recall:} say the Japanese for selfish, then the Japanese for to agree, then say \emph{hantai suru} again
\end{itemize}

\begin{tcolorbox}[breakable,colback=teal!4,colframe=teal!35!black,title={Before you move on}]
What does \ja{はんたいする} mean? (\textquotedblleft{}To be against; to oppose\textquotedblright{}.) Say it once more.
\end{tcolorbox}
\section[\texorpdfstring{\ja{ぎろん} (giron)}{ぎろん (giron)}]{\ja{ぎろん} (giron) — a discussion; an argument}

\label{lesson:JA-C317-giron}



\begin{quote}
Before the new one: say the Japanese for to agree, then the Japanese for to be against.
\end{quote}

\subsection*{You'll want to know: \ja{ぎろん}}
\textbf{\ja{ぎろん}} — \emph{giron} — \textquotedblleft{}a discussion; an argument\textquotedblright{}.

\begin{quote}
\textbf{\ja{ぎろんが} \ja{ながく} \ja{なったので、つかれました。}} — \emph{giron ga nagaku natta node, tsukaremashita} — \textquotedblleft{}The discussion dragged on, so I got tired.\textquotedblright{}
\end{quote}

\begin{cousinweb}
\textbf{\ja{ぎろん}} (\emph{giron}) is Sino-Japanese (kango): two Chinese characters read with their on'yomi, readings borrowed from Chinese. They are \emph{gi}, \textquotedblleft{}debate\textquotedblright{} (Middle Chinese \emph{ngjeH}) and \emph{ron}, \textquotedblleft{}argument\textquotedblright{} (\emph{lwon}) (KANJIDIC2; Middle Chinese in Baxter's transcription).
\end{cousinweb}

\begin{grammarlens}[title={What you've built}]
One more word for giving reasons and opinions.
\end{grammarlens}

\subsection*{Your turn}
\begin{itemize}
\raggedright
  \item \emph{Say it:} \emph{giron}
  \item \emph{Say it:} \emph{giron}, once more
  \item \emph{Say it:} \emph{giron ga nagaku natta node, tsukaremashita} — The discussion dragged on, so I got tired.
  \item \emph{Recall:} say the Japanese for to agree, then the Japanese for to be against, then say \emph{giron} again
\end{itemize}

\begin{tcolorbox}[breakable,colback=teal!4,colframe=teal!35!black,title={Before you move on}]
What does \ja{ぎろん} mean? (\textquotedblleft{}A discussion; an argument\textquotedblright{}.) Say it once more.
\end{tcolorbox}
\section[\texorpdfstring{\ja{ていあんする} (teian suru)}{ていあんする (teian suru)}]{\ja{ていあんする} (teian suru) — to propose; to suggest a plan}

\label{lesson:JA-C317-teian-suru}



\begin{quote}
Before the new one: say the Japanese for to be against, then the Japanese for a discussion.
\end{quote}

\subsection*{You'll want to know: \ja{ていあんする}}
\textbf{\ja{ていあんする}} — \emph{teian suru} — \textquotedblleft{}to propose; to suggest a plan\textquotedblright{}.

\begin{quote}
\textbf{\ja{やすいので、この} \ja{ほうほうを} \ja{ていあんします。}} — \emph{yasui node, kono hōhō o teian shimasu} — \textquotedblleft{}It is cheap, so I propose this method.\textquotedblright{}
\end{quote}

Polite \textbf{\ja{ていあんします}}, \emph{teianshimasu}; a \ja{する} verb. Past (polite): \textbf{\ja{じょうしが} \ja{ていあんしました。}}, \emph{jōshi ga teian shimashita}.

\begin{cousinweb}
\textbf{\ja{ていあんする}} (\emph{teian suru}) is built from a Sino-Japanese noun and the native verb \emph{suru}, \textquotedblleft{}to do\textquotedblright{}. The noun, \emph{teian}, is two Chinese characters read with their on'yomi, readings borrowed from Chinese: \emph{tei}, \textquotedblleft{}propose\textquotedblright{} and \emph{an}, \textquotedblleft{}plan\textquotedblright{} (Middle Chinese \emph{'anH}) (KANJIDIC2; JMdict; Middle Chinese in Baxter's transcription).
\end{cousinweb}

\begin{grammarlens}[title={What you've built}]
One more word for giving reasons and opinions.
\end{grammarlens}

\subsection*{Your turn}
\begin{itemize}
\raggedright
  \item \emph{Say it:} \emph{teian suru}
  \item \emph{Say it:} \emph{teian suru}, once more
  \item \emph{Say it:} \emph{yasui node, kono hōhō o teian shimasu} — It is cheap, so I propose this method.
  \item \emph{Recall:} say the Japanese for to be against, then the Japanese for a discussion, then say \emph{teian suru} again
\end{itemize}

\begin{tcolorbox}[breakable,colback=teal!4,colframe=teal!35!black,title={Before you move on}]
What does \ja{ていあんする} mean? (\textquotedblleft{}To propose; to suggest a plan\textquotedblright{}.) Say it once more.
\end{tcolorbox}
\section[\texorpdfstring{\ja{けつろん} (ketsuron)}{けつろん (ketsuron)}]{\ja{けつろん} (ketsuron) — a conclusion}

\label{lesson:JA-C317-ketsuron}



\begin{quote}
Before the new one: say the Japanese for a discussion, then the Japanese for to propose.
\end{quote}

\subsection*{You'll want to know: \ja{けつろん}}
\textbf{\ja{けつろん}} — \emph{ketsuron} — \textquotedblleft{}a conclusion\textquotedblright{}.

\begin{quote}
\textbf{\ja{けつろんは} \ja{あしたの} \ja{かいぎで} \ja{きめた} \ja{ほうが} \ja{いいです。}} — \emph{ketsuron wa ashita no kaigi de kimeta hō ga ii desu} — \textquotedblleft{}We had better reach a conclusion at tomorrow's meeting.\textquotedblright{}
\end{quote}

\begin{cousinweb}
\textbf{\ja{けつろん}} (\emph{ketsuron}) is Sino-Japanese (kango): two Chinese characters read with their on'yomi, readings borrowed from Chinese. They are \emph{ketsu}, \textquotedblleft{}tie\textquotedblright{} (Middle Chinese \emph{ket}) and \emph{ron}, \textquotedblleft{}argument\textquotedblright{} (\emph{lwon}) (KANJIDIC2; Middle Chinese in Baxter's transcription).
\end{cousinweb}

\begin{grammarlens}[title={What you've built}]
One more word for giving reasons and opinions.
\end{grammarlens}

\subsection*{Your turn}
\begin{itemize}
\raggedright
  \item \emph{Say it:} \emph{ketsuron}
  \item \emph{Say it:} \emph{ketsuron}, once more
  \item \emph{Say it:} \emph{ketsuron wa ashita no kaigi de kimeta hō ga ii desu} — We had better reach a conclusion at tomorrow's meeting.
  \item \emph{Recall:} say the Japanese for a discussion, then the Japanese for to propose, then say \emph{ketsuron} again
\end{itemize}

\begin{tcolorbox}[breakable,colback=teal!4,colframe=teal!35!black,title={Before you move on}]
What does \ja{けつろん} mean? (\textquotedblleft{}A conclusion\textquotedblright{}.) Say it once more.
\end{tcolorbox}
